\chapter{Related Work}\label{ch:related_work}

\Cref{ch:introduction} identified a gap between style authoring and tile data preparation, two concerns that remain independent in all existing open-source tooling.
The literature relevant to close that gap is imbalanced.
Tile data pruning has been attacked head-on by several open-source tools, each addressing a slice of the problem and no one uniting this work with style rewriting.
Tile encoding has seen substantial recent activity around columnar formats, but the style is always treated as a downstream consumer rather than an input.
Style-level optimization has attracted comparatively little research attention, and the closest precedents come entirely from outside the map domain.
Benchmarking methodology, finally, is the area where the prior art is mature, and we mostly borrow rather than contribute.

We work through these areas in the same order:
\cref{sec:rw_tile_data} for tile pruning and schema design,
\cref{sec:rw_tile_encoding} for encoding and columnar compression,
\cref{sec:rw_style_optimizations} for style-level work and its \ac{CSS}-minification cousins,
and \cref{sec:rw_benchmarking} for rendering benchmarking.

\section{Tile Data Optimization}\label{sec:rw_tile_data}

The most direct approach to reduce tile transfer cost is to remove unreferenced data from a style.

Mapbox developed \texttt{vtshaver}\footnote{\url{https://github.com/mapbox/vtshaver}, last accessed 2025-10-20}, a C++/Node.js tool that takes a vector tile and a Mapbox GL style, analyses which layers, features, and properties the style references, and strips everything else.
This reduces tile size and decoding cost, but \texttt{vtshaver} operates exclusively on the data side.
It does not modify the style itself and therefore cannot simplify filter expressions, merge redundant layers, or reduce per-frame rendering overhead.
It also only shaves on a zoom-property-unaware row-column approach.

Mapbox also offers \emph{style-optimized vector tiles} as a proprietary server-side feature~\cite{StyleoptimizedVectorTiles}.
When a tile request includes style metadata, the server analyses the style's sources, filters, and zoom ranges, then strips all data not visible for that style.
This is the closest commercial precedent for style-data co-optimizations, but it is opaque, closed-source, and \ac{API}-only.
It cannot be extended, audited, or deployed independently.
The exact set of optimizations is not disclosed and given that MapBox does not have a open research stance (their terms of service explicitly forbid this), reverse-engineering would likely be illegal\footnote{\url{https://media.ccc.de/v/36c3-11089-reflections_on_the_new_reverse_engineering_law}, last accessed 2026-05-21}.

Vilardaga built \texttt{vt-optimizer}\footnote{\url{https://github.com/ibesora/vt-optimizer}, last accessed 2025-10-20}, which drops invisible or unused layers and fields in both the style and data.
The tool also provides an inspection mode that checks conformance to Mapbox's recommended tile size limits (\qty{50}{\kilo\byte} average, \qty{500}{\kilo\byte} maximum)\footnote{\url{https://docs.mapbox.com/mapbox-tiling-service/guides/frequentlyaskedquestions/}, last accessed 2026-05-23}.
However, the optimizations are limited to visibility-based pruning.
It does not perform expression-level rewrites, cross-layer merging, or data-driven constant folding.

Tippecanoe\footnote{\url{https://github.com/felt/tippecanoe}, last accessed 2026-04-21} is the standard open-source tool for building vector tilesets from large \ac{GeoJSON} or FlatGeobuf\footnote{\url{http://flatgeobuf.org/}, last accessed 2026-05-11} collections.
It implements several data-level optimizations heuristics: feature dropping based on density budgets (\texttt{--drop-densest-as-needed}), polygon coalescing at low zoom levels, Visvalingam line simplification~\cite{visvalingamLineGeneralisationRepeated1993}, and attribute exclusion.
The line simplification algorithms used throughout the geospatial toolchain (Douglas-Peucker~\cite{douglasALGORITHMSREDUCTIONNUMBER1973} for its computational simplicity and Visvalingam-Whyatt~\cite{visvalingamLineGeneralisationRepeated1993} for its area-based perceptual fidelity) are instances of \emph{cartographic generalization}, the principled reduction of map detail to match the target scale.
Tile data at different zoom levels is, in effect, a discrete multi-scale generalization hierarchy.
However, this thesis does not contribute to new generalization algorithms.
These heuristics are effective but purely data-driven.
Tippecanoe has no knowledge of which style will render the tiles and therefore cannot tailor its decisions to the rendering workload.

From a cartographic perspective, we understand style-driven tile shaving as a form of \emph{model generalization}~\cite{mcmasterGeneralizationDigitalCartography1992}.
The rendering style serves as the generalization specification, and the pruning advisory selects which data survives the reduction, analogous to how a target scale determines which features are retained in traditional multi-scale generalization.
The distinction is that classical generalization selects features by geometric criteria (area, length, density), whereas style-driven shaving selects by attribute criteria (properties and values referenced by the style).

At a higher level of abstraction, tile schemas themselves embody data-level design decisions.
The OpenMapTiles schema~\cite{VectorTilesAre2025} defines which attributes are included at which zoom levels for OpenStreetMap data.
Wallner et al.~\cite{wallnerOpenSourceVector2022} describe a PostGIS-based pipeline for dynamic vector tile creation with runtime geometry simplification and \ac{MVT} encoding for spatial data infrastructures.
Tile generators such as Planetiler\footnote{\url{https://github.com/onthegomap/planetiler}, last accessed 2026-04-21} and Tilemaker\footnote{\url{https://tilemaker.org/}, last accessed 2026-04-21} apply schema-aware rules during the initial tile build step, and Protomaps basemap\footnote{\url{https://protomaps.com}, last accessed 2026-04-21} ships a style-specific schema in which column selection is effectively hard-coded into the tile generator.
These systems define the data that enters the pipeline.
Our optimizations operate downstream, on the data and style that emerge from such pipelines, and do not require re-running the generator.

A rendering style is a declarative specification over a tabular feature set, which puts style-driven tile pruning in the same shape as the column- and predicate-skipping problems that database query optimizers have long attacked.
Rule-based query rewriting was established as a distinct optimization phase by Pirahesh et~al.'s Starburst system~\cite{piraheshExtensibleRuleBased1992}, and projection and predicate pushdown~\cite{graefeQueryEvaluationTechniques1993} are the canonical mechanisms by which modern engines such as Spark SQL~\cite{armbrustSparkSQLRelational2015} avoid materialising columns the query never reads.
Columnar storage systems built on Parquet (whose nested representation originates in Dremel's record shredding and assembly~\cite{melnikDremelInteractiveAnalysis}) or ORC depend on this pushdown, and Abadi et~al.~\cite{abadiMaterializationStrategiesColumnOriented2007} characterised the early versus late materialisation trade-off it implies.
Our data level optimisations are explained in \cref{sub:data_optimizations}.

\section{Tile Encoding and Format Innovation}\label{sec:rw_tile_encoding}

The dominant wire format for vector tiles is \ac{MVT}\footnote{\url{https://github.com/mapbox/vector-tile-spec/tree/master/2.1}, last accessed 2026-04-21}, which uses Protocol Buffers\footnote{\url{https://protobuf.dev/}, last accessed 2026-04-21} to encode features in a row-oriented layout.
\ac{MVT} has been adopted as a conformance class in the \ac{OGC} API-Tiles standard~\cite{OGCAPITiles}, making it the baseline format for any standards-compliant tile delivery.
Its row orientation means that accessing a single property column requires scanning the entire feature table, and its protobuf encoding does not exploit the statistical properties of individual columns.

Tremmel~\&~Zink.~\cite{tremmelMapLibreTileNext2025b} introduced the MapLibre Tile (MLT) format, a columnar alternative that achieves up to $3\times$ tile size reduction on encoded tilesets and up to $6\times$ on certain large tiles relative to \ac{MVT}.
\ac{MLT} uses column-oriented storage with lightweight compression (delta encoding, run-length encoding, dictionary encoding, and \ac{FSST}~\cite{bonczFSSTFastRandom2020}) and is designed for \ac{GPU}-accelerated processing.
At the archive layer, Tremmel~\cite{tremmelCOMTILESCASESTUDY2023} proposed COMTiles, a cloud-optimized tile archive format that complements the per-tile encoding of \ac{MLT} by addressing the container and delivery layer for planet-scale tilesets.
However, their work does not consider styles.
The encoding strategy is chosen without knowledge which properties the renderer will actually access, leaving potential for further size reduction through style-aware pruning before encoding.
Removing properties and features changes column cardinalities and value distributions.
For instance, pruning rare road classes can convert a high-cardinality string column into one amenable to dictionary or run-length encoding, so the optimal encoding strategy for the shaved tile differs from the original.

The encoding techniques used in \ac{MLT} draw on a rich literature in columnar database compression.
Lemire and Boytsov~\cite{lemireDecodingBillionsIntegers2015} introduced \ac{SIMD}-vectorised integer compression \ac{FastPFOR} achieving billions of integers per second in decompression throughput.
Boncz et al.~\cite{bonczFSSTFastRandom2020} developed \ac{FSST}, a lightweight string compression scheme using a static symbol table that supports random access without decompressing surrounding data.
Kuschewski et al.~\cite{kuschewskiBtrBlocksEfficientColumnar2023} demonstrated with BtrBlocks that cascading multiple encoding schemes (including \ac{FSST}, bit-packing, \ac{RLE}, and dictionary encoding) with automatic per-block strategy selection achieves 2.2$\times$ faster scans than Parquet.
The general finding that no single encoding is universally optimal, and that data-driven strategy selection is necessary, motivates the multi-strategy encoding competition we develop in this thesis.

The problem of automatic encoding selection in columnar databases has a substantial lineage.
Abadi et~al.~\cite{abadiIntegratingCompressionExecution2006} introduced decision rules for choosing among \ac{RLE}, dictionary, and bit-packing based on column properties (cardinality, sort order, run lengths) in the C-Store system.
Damme et~al.~\cite{dammeComprehensiveExperimentalSurvey2019} evaluated dozens of lightweight integer encodings and developed a cost-based selection strategy.
Boissier~\cite{ganBaguaScalingDistributed2021} formulated encoding selection as a budget-constrained optimization problem.
And Zeng et~al.~\cite{zengEmpiricalEvaluationColumnar2023} empirically evaluated Parquet and ORC internals, confirming that dictionary encoding should be a default candidate and that decompression speed matters more than compression ratio for integer encodings.

Our encoding competition differs from these systems by operating at tile granularity (\qtyrange{10}{100}{\kilo\byte}) rather than on large analytical datasets.
We also include application-specific optimizations like spatial sort-order dimension (ID, Hilbert, Morton) that database encoders do not consider.
This competition is an instance of \emph{instance-optimized} system design as proposed by Kraska~\cite{dingSageDBInstanceOptimizedData2022}: rather than committing to a fixed encoding plan, the encoder uses the data distribution of each tile to select optimal physical representations, exactly the approach advocated by the SageDB vision~\cite{dingSageDBInstanceOptimizedData2022} and realised concretely by Ding et~al.~\cite{dingInstanceOptimizedDataLayouts2021}, who showed that instance-optimized data layouts reduce blocks accessed by 93\% on cloud analytics workloads, and at the column-block level by BtrBlocks~\cite{kuschewskiBtrBlocksEfficientColumnar2023}.
Alternative approaches to encoding selection include Fehér et~al.'s~\cite{feherAdaptiveColumnCompression2022} adaptive compression family, which selects encodings at runtime using segment-level measurements, and Cen et~al.'s~\cite{cenLEALearnedEncoding2021} LEA, a learned encoding advisor that uses ML models to predict the best encoding per column, achieving 19\% lower latency and 26\% less space than heuristic selection on TPC-H.
Our trial-based competition is closest in spirit to the instance-optimized approach, but substitutes exhaustive evaluation over a bounded candidate space for the learned or heuristic strategies used in database systems.

Gaffuri~\cite{gaffuriWebMappingVector2012} presented one of the earliest academic treatments of vector tile architecture for web maps, proposing a framework integrating tiling, spatial indexing, and generalization for multi-scale data.
His work established the architectural vocabulary that \ac{MVT} and \ac{MLT} build on, but predates the shift to columnar encoding and does not consider style-awareness as a dimension of the encoding problem.
That is precisely the dimension we add to the tile encoding pipeline.

GeoParquet uses the same per-page encoding competition (dictionary, \ac{RLE}, delta, bit-packing) that \ac{MLT} applies per stream, and at first glance it is the obvious off-the-shelf target for vector tile delivery.
The reason it is not, in practice, is granularity.
Parquet's per-column-chunk metadata is amortised over the large row groups of analytical workloads, but at the \qtyrange{10}{100}{\kilo\byte} tile sizes we care about that overhead dominates the payload.
A purpose-built tile format earns its keep here despite the ecosystem support that adopting GeoParquet would have brought along.
Our encoding level optimisations are explained in \cref{sub:encoding_strategies}.

\section{Style Optimization}\label{sec:rw_style_optimizations}

In contrast to data-level optimizations, style-level optimization techniques have attracted comparatively little research attention.

\begin{sloppypar}
The MapLibre style specification includes the utilities \texttt{gl\allowbreak-style\allowbreak-validate} and \texttt{gl\allowbreak-style\allowbreak-migrate}\footnote{\url{https://www.maplibre.org/maplibre-style-spec/}, last accessed 2025-10-24}, which validate style documents against the specification and mechanically migrate legacy function syntax to the expression language.
These tools are not simplifiers but deliberately conservative transformers, and the mechanical migration they produce often contains redundancy that a simplifier could improve.
Deeply nested \mlop{case} chains where a \mlop{match} is sufficent, explicit \mlop{typeof} guards that are trivially true after the migration, and boolean expressions that could be flattened.
A simplification pass is therefore a natural complement to migration, not a replacement for it.
\end{sloppypar}

More broadly, style documents in practice accumulate redundancy through several channels: iterative authoring that leaves superseded properties in place, copy-paste development across layers, automated migration that introduces mechanically correct but verbose expression forms, and multi-author editing where conventions drift over successive contributions.
The result is that production styles routinely contain default-valued properties, unreachable filter branches, and duplicated paint blocks, creating a practical demand for automatic simplification that no existing tool meets.

Stamen Design published \texttt{mapbox-gl-style-remove-defaults}\footnote{\url{https://github.com/stamen/mapbox-gl-style-remove-defaults}, last accessed 2026-5-27}, a tool that removes properties set to their specification default values from Mapbox GL style \ac{JSON}.
This corresponds to the default stripping pass we describe in \cref{ch:methods}, implemented as a standalone tool.
This ecosystem tooling demonstrates that individual style optimizations are recognized as valuable by practitioners, but no existing tool encomposes multiple optimizations passes.

The closest academic analogue lies outside the map domain entirely.
Hague et al.~\cite{hagueCSSMinificationConstraint2019} formalise \ac{CSS} rule merging as a Max-SAT problem over a node-weighted bipartite graph of rules and properties with ordering constraints from the cascade, and their SatCSS tool beats six production \ac{CSS} minifiers on real benchmarks.
Our layer-merging pass (\cref{sub:structural_optimizations}) attacks the same declarative-rule-merging-under-ordering-constraints problem, but uses heuristics instead.

More broadly, rule-based optimization is a well-studied paradigm in databases and compilers, and a rule-based style rewriter sits between two architectural reference points.
Pirahesh et~al.~\cite{piraheshExtensibleRuleBased1992} established \emph{rule-based query rewriting} as a distinct optimization phase in the Starburst system, in which declarative transformation rules are applied in priority order before cost-based plan selection.
The modern state of the art for avoiding rule-ordering sensitivity is \emph{equality saturation}~\cite{tateEqualitySaturationNew2009,willseyEggFastExtensible2021}, in which an e-graph simultaneously represents all equivalent forms and selects an optimum after saturation.
Our style level optimisations are explained in \cref{sub:expression_optimizations,sub:structural_optimizations}.

\section{Rendering Performance Benchmarking}\label{sec:rw_benchmarking}

Several studies have benchmarked vector map rendering performance.

Netek et al.~\cite{netekPerformanceTestingVector2020} compared vector and raster tile performance across eight scenarios, measuring loading time, data transfer size, and \ac{HTTP} request count.
They find that vector tiles offer better loading times, though raster tiles may be preferable on very slow connections.

Balla and Gede~\cite{ballaVectorDataRendering2025} benchmarked \ac{GeoJSON} rendering performance across Leaflet, OpenLayers, MapLibre GL~JS, and Mapbox GL~JS with varying feature counts.
Below 10\,000 features, \ac{SVG}-based renderers (Leaflet, OpenLayers) are faster.
Above 50\,000 features, \ac{WebGL}-based renderers (Mapbox GL~JS, MapLibre GL~JS) dominate due to \ac{GPU} acceleration.

\begin{sloppypar}
Rechsteiner~\cite{rechsteinerPerformancevergleichOpenSourceVectorTilesServerloesungenZur2024} compared open-source vector tile server implementations.
Since the systems under test offer different feature sets, the evaluation partly compares incomparable aspects and does not address client-side rendering optimizations or style-level improvements.
\end{sloppypar}

Georges et~al.~\cite{georgesStatisticallyRigorousJava2007} and Hoefler and Belli~\cite{hoeflerScientificBenchmarkingParallel2015} established the methodological foundations for statistically rigorous performance evaluation.
The specific techniques drawn from their work (bootstrap confidence intervals, warm-up handling, and metric separation) are introduced in \cref{sec:benchmarking}.

These benchmarking studies establish the measurement landscape but focus on comparing \emph{systems} rather than evaluating the impact of \emph{optimizations passes} applied to a fixed system.
We complement them by introducing a cumulative ablation methodology that isolates the contribution of each optimizations pass within a single renderer (MapLibre GL~JS).

\section{Positioning}\label{sec:rw_positioning}

Each strand of this literature solves part of the problem in isolation.
The data-side tools (\texttt{vtshaver}, \texttt{vt-optimizer}, Tippecanoe) prune tile content but never read or rewrite the style document that drives the renderer.
Mapbox's closed-source \enquote{style-optimized tiles} service is the only precedent that genuinely couples style and data, and its terms of service explicitly preclude the independent evaluation needed to confirm what it actually does.
The style-side work is sparser still.
Stamen's default-stripping utility and a handful of validators cover individual optimizations, but no published tool composes multiple passes, and the closest academic analogue (\ac{CSS} minification) lives in a different domain altogether.
Tremmel and Zink's \ac{MLT} format advances the encoding layer but picks strategies without consulting the style.
The benchmarking literature is in better shape, and we adopt its statistical methodology rather than extending it.

\begin{chaptersummary}
No prior work composes style and data optimizations into a single open pipeline (the data-side tools ignore the style, the only commercial precedent is closed-source, and the encoding format's strategy selection is style-blind).
Our contribution is the joint $(S, T) \to (S', T')$ optimization formalised in \cref{sec:problem_statement}, together with a cumulative ablation that distinguishes additive from synergistic composition.
\end{chaptersummary}
